\section{Posentrenamiento, RLHF y Machines of Loving Grace}

\subsection{Flujo moderno de posentrenamiento}

El posentrenamiento moderno incluye varias etapas: fine-tuning supervisado (SFT), RLHF, Constitutional AI (RLAIF), generación de datos sintéticos. Es difícil atribuir la mejora del modelo al preentrenamiento o al posentrenamiento: distintos equipos avanzan cada uno en un tramo distinto de la ``carrera de relevos''.

La ventaja de Anthropic en RL \textit{``tal vez sea la mejor''}, pero Dario insiste en que esa ventaja normalmente no viene de un método mágico secreto: \textit{``es una cuestión de práctica aburrida y oficio''}. Es como diseñar un avión: lo que determina la calidad no es un solo avance, sino ``el oficio cultural de cómo pensamos el proceso de diseño''.

\begin{importantbox}{La esencia de RLHF: liberar, no crear}
\textit{``RLHF no hace más inteligente al modelo. Tampoco se trata solo de que el modelo parezca más inteligente. RLHF tiende un puente entre los humanos y el modelo.''}

Es como personas muy inteligentes pero poco hábiles para expresarse: RLHF ``libera'' (un-hobble) las capacidades que el modelo ya posee. Tomando el término de Leopold Aschenbrenner: el preentrenamiento le da al modelo el potencial de conocimiento y razonamiento, y RLHF hace que ese potencial se manifieste en una forma que los humanos puedan usar.

Sin embargo, el RL sí tiene el potencial de hacer al modelo genuinamente más inteligente y con mejor razonamiento, pero el RLHF actual sigue siendo sobre todo ``liberación'' y no ``mejora''. El preentrenamiento todavía representa la mayor parte del costo de entrenamiento, pero esa proporción podría invertirse en el futuro.
\end{importantbox}

\subsection{Constitutional AI}

El artículo de diciembre de 2022 tiene una idea central: la IA juzga, con base en un conjunto de principios legibles por humanos, cuál respuesta es mejor, formando un ciclo triangular de autojuego. En la práctica, CAI se combina con RLHF y otros métodos: \textit{``una herramienta más en la caja de herramientas''}.

La redacción de la constitución: los principios básicos son casi un consenso universal (prohibir CBRN, el Estado de derecho democrático básico). Más allá de eso, el modelo debe ser \textit{``más neutral, sin inclinarse hacia una postura específica''} —como un ``consejero sabio'' y no un predicador.

El Model Spec de OpenAI tomó una dirección similar—esto es precisamente una muestra de la Race to the Top. John Schulman (ahora en Anthropic) participó en la redacción del Model Spec. Es posible que Anthropic también publique su propio Model Spec.

\subsection{Machines of Loving Grace: por qué Dario escribió este ensayo}

Dario se ha enfocado durante mucho tiempo en los riesgos de la IA, pero se dio cuenta de que hablar solo de riesgos hace que la mente solo piense en riesgos. \textit{``La razón por la que intentamos prevenir estos riesgos no es porque le tengamos miedo a la tecnología.''}

Si logramos ``atravesar con éxito el campo minado...del otro lado están todas estas cosas maravillosas''. Él quería que alguien del bando de los riesgos expusiera en serio los beneficios de la IA.

\begin{importantbox}{Más allá del catastrofismo y el aceleracionismo}
\textit{``Esto no es catastrofistas contra aceleracionistas...si de verdad entiendes hacia dónde va la IA, vas a apreciar genuinamente sus beneficios y también vas a tomarte muy en serio cualquier riesgo que pueda arruinarla.''}

Crítica al optimismo tecnológico vago: \textit{``la ciudad reluciente...acelerar, acelerar y acelerar más...pero ¿de qué, exactamente, estás entusiasmado?''} Lo que Dario exige son beneficios concretos y demostrables, no un optimismo vacío.
\end{importantbox}

\subsection{La definición de ``IA poderosa'' y su cronograma}

A Dario no le gusta el término ``AGI'': no hay un umbral discreto, solo un crecimiento exponencial uniforme. Él define la ``IA poderosa'' como:

\begin{enumerate}
    \item Más inteligente que un premio Nobel (en su mejor momento) en la mayoría de las disciplinas relevantes
    \item Capaz de usar todas las modalidades (texto, imágenes, código, herramientas)
    \item Capaz de ejecutar tareas de forma independiente durante horas, días o incluso semanas
    \item No necesariamente tiene cuerpo, pero puede controlar herramientas físicas (robots, equipo de laboratorio)
    \item Los recursos usados para entrenarla pueden redirigirse a ejecutar \textbf{millones de copias}
    \item Cada copia trabaja de 10 a 100 veces más rápido que un humano
\end{enumerate}

Cronograma: \textit{``si extrapolas esta línea...2026 o 2027. Lo más probable es que haya algún pequeño retraso.''} Pero \textit{``el número de mundos en los que esto no ocurre se está reduciendo rápidamente.''}

\subsection{La singularidad frente al estancamiento: dos extremos equivocados}

Dario refuta dos posturas extremas:

\begin{warningbox}{Ambos están equivocados}
\textbf{Extremo uno: la singularidad}—la IA construye una IA más rápida, hay una mejora recursiva, y una IA sobrehumana aparece y cinco días después todo ha sido inventado. Por qué está equivocado: las leyes de la física (el hardware tarda en fabricarse), la complejidad (algunos sistemas requieren experimentos, no modelos) y una impredecibilidad similar a la del problema de los tres cuerpos. \textit{``con los sistemas biológicos...siempre es mejor ejecutar el experimento directamente que modelarlo, sin importar qué tan bueno sea quien modela.''}

\textbf{Extremo dos: el estancamiento}—al igual que con la revolución de las computadoras, la mejora en la productividad resulta decepcionante. La célebre frase de Robert Solow: \textit{``puedes ver la revolución de las computadoras en todas partes, menos en las estadísticas de productividad.''} Tyler Cowen cree que el cambio tomará entre 50 y 100 años. Por qué está equivocado: la escala de tiempo es demasiado larga.
\end{warningbox}

La postura intermedia de Dario: \textbf{de 5 a 10 años}, ni 50 a 100 años ni 5 a 10 horas. \textit{``los obstáculos se irán desmoronando poco a poco, y luego, de repente, se derrumbarán todos.''} La inercia de las instituciones humanas es el verdadero cuello de botella.

\subsubsection{Cómo ocurre el cambio en las grandes organizaciones}

Dos fuerzas impulsan el cambio:
\begin{enumerate}
    \item Una minoría de visionarios dentro de la organización, que ven el panorama completo
    \item El fantasma de la competencia—\textit{``miren, esa gente ya lo está haciendo, nos van a comer el almuerzo''}
\end{enumerate}

Los visionarios por sí solos no bastan, pero la competencia les da \textit{``viento a favor''}. Esto se repite igual en la banca, en los gobiernos, dentro del propio gobierno de Estados Unidos. \textit{``la inercia es enormemente poderosa, pero al final la innovación encuentra la forma de abrirse paso.''} Igual que con la difusión de la Scaling Hypothesis: \textit{``sentíamos que teníamos un secreto que casi nadie conocía, y unos años después todos lo sabían.''}

\subsection{El impacto de la IA en la programación}

La programación será el ámbito que cambie \textbf{más rápido}, por dos razones:
\begin{enumerate}
    \item Es lo más cercano a quienes construyen la IA—la ``distancia'' a los desarrolladores de IA equivale al tiempo que falta para ser disruptido
    \item El modelo puede cerrar el ciclo: escribir código, ejecutarlo, ver los resultados e iterar
\end{enumerate}

SWE-bench: pasó de 3\% a 50\% en 10 meses. Dario predice que en otros 10 meses se acercará a 90\%.

\begin{knowledgebox}{La ventaja comparativa y la analogía del procesador de textos}
Cuando la IA hace el 80\% del trabajo de codificación, el 20\% restante que le queda al humano (el diseño de alto nivel, la arquitectura, la UX) se vuelve, en cambio, más apalancado. Igual que el procesador de textos hizo instantánea la escritura y la maquetación, y toda la atención se desplazó hacia las \textbf{ideas}.

\textit{``durante más tiempo del que esperamos...la pequeña parte del trabajo que los humanos todavía hacen se expandirá hasta llenar todo su trabajo.''} Al final, la IA superará a los humanos en todos los aspectos—\textit{``y en ese momento la humanidad tendrá que pensar colectivamente cómo responder.''} Pero en el corto y mediano plazo (de 2 a 4 años): \textit{``la programación como trabajo no va a cambiar, pero la naturaleza de la programación sí va a cambiar.''}
\end{knowledgebox}

\subsubsection{El futuro de los IDE y las herramientas de desarrollo}

\textit{``incluso si la calidad del modelo dejara de mejorar, simplemente aumentar la productividad de las personas ya sería una oportunidad enorme.''} Un IDE puede encargarse de: análisis estático, detección de bugs, organización del código, cobertura de pruebas, además de que la IA escriba y ejecute código.

La estrategia de Anthropic: \textit{``dejar que florezcan cien flores''} —empoderar a empresas como Cursor, Cognition y Expo que construyen sobre la API, en lugar de competir con sus clientes (al menos por ahora).

\begin{importantbox}{El potencial de la IA en la biología—el ``siglo XXI comprimido''}
La IA poderosa podría comprimir el avance de la biología previsto entre 2027 y 2100 en el período de 2027 a 2032—curar la mayoría de los cánceres, prevenir enfermedades infecciosas y duplicar la esperanza de vida humana.

El día a día del biólogo del futuro: como tener 1,000 estudiantes de posgrado de IA más inteligentes que uno mismo—capaces de revisar la bibliografía, encargar equipo, hacer experimentos, examinar imágenes y escribir código estadístico. \textit{``antes un profesor tenía 50 estudiantes de posgrado, ahora tienes 1,000, y son más inteligentes que tú.''}

Al final los papeles se invertirán: el sistema de IA se convertirá en el investigador principal (PI), y dirigirá a humanos y a otras IA.
\end{importantbox}

\subsection{El sentido, el trabajo y la concentración de poder}

El ensayo ``Machines of Loving Grace'' de Dario se planeó originalmente con 2 o 3 páginas, y terminó expandiéndose a 40 o 50. Sobre el sentido, propone un experimento mental: si vivieras 60 años dentro de una simulación y luego descubrieras que era un juego, ¿se te habría arrebatado el sentido? Él cree que no: \textit{``lo que importa es el proceso...eso es lo que muestra el tipo de persona que eres.''}

Si se maneja mal, un mundo con IA podría carecer de sentido—pero eso es una decisión de diseño social, no algo inevitable. Más importante aún: la mayoría de la gente en el mundo pasa su tiempo ``luchando por sobrevivir'', y su vida mejorará enormemente gracias a la IA. \textit{``tratar el sentido como lo único importante es, en cierta forma, un privilegio de una minoría con suerte económica.''}

\begin{warningbox}{La mayor preocupación de Dario: la concentración de poder}
No es la pérdida de sentido ni el desempleo, sino \textbf{la concentración y el abuso del poder}.

\textit{``la IA aumenta la cantidad total de poder que hay en el mundo, y si concentras ese poder y abusas de él, el daño que puedes causar es incalculable.''}

Estructuras como la dictadura y el autoritarismo—en las que una minoría explota a la mayoría—son el futuro de la IA que más le preocupa a Dario. Esto es más fundamental que los problemas técnicos de seguridad de la IA, y también más difícil de resolver. Los problemas técnicos tienen soluciones técnicas, pero la distribución del poder es un problema político.
\end{warningbox}

\subsection{Resumen de la sección}

El núcleo del posentrenamiento es un ``oficio cultural de práctica'' y no un algoritmo secreto. RLHF libera las capacidades del modelo en lugar de crearlas. La IA poderosa podría llegar entre 2026 y 2027, y transformará por completo la biología y la programación. El mayor riesgo no son los problemas técnicos de la IA en sí, sino cómo se distribuye el poder.

%% ==================== Part II: Amanda Askell ====================
\part{Amanda Askell: la construcción de la personalidad de Claude}

